\section{The Chow Groups}

\bdefn

    Let $X$ be any algebraic variety.
    A {\emphcolor $k$-cycle} on $X$ is a formal sum of subvarieties (closed integral subschemes) of $X$
    of dimension $k$.
    In other words the group of $k$-cycles on $X$ is the free Abelian group on $X$'s $k$-dimensional
    subvarieties, and is denoted $Z_k(X)$.
    $X$'s {\emphcolor group of cycles} is the graded Abelian group $Z(X)=\bigoplus_kZ_k(X)$.

    A cycle $\sum n_iY_i\in Z(X)$ is {\emphcolor effective} if all of the $n_i$ are nonnegative.
    A {\emphcolor divisor} is an $(n-1)$-cycle (where $n=\dim X$).

\edefn

For a variety $X$ let $K(X)$ denote its field of rational functions.
It is the fraction field of any one of its stalks, and is equal to its stalk at its generic point.
For a subvariety $Y\subseteq X$, we denote by $\O_{X,Y}$ the stalk of $X$ at $Y$'s generic point.

\bdefn

    Let $Z\subseteq X$ be a variety of codimension $1$, then for $f\in K(X)^\times$ we denote its
    {\emphcolor order of vanishing} along $Z$ as follows.
    Because $\O_{X,Z}$ is a one-dimensional local domain, for any nonzero $f\in\O_{X,Z}$ the quotient ring
    $\O_{X,Z}/(f)$ is a zero-dimensional Noetherian ring, and thus has finite length.
    We define $f$'s order of vanishing to be
    $$ \ord_Z(f) = \length_{\O_{X,Z}}\O_{X,Z}/(f) . $$
    Any rational function $r\in K(X)^\times$ has the form $f/g$ for $f,g\in\O_{X,Z}$ and so we define
    $$ \ord_Z(r) = \ord_Z(f) - \ord_Z(g) . $$

\edefn

It remains to show that the order of vanishing is well-defined.
To do this we note that $\ord_Z(fg)=\ord_Z(f)+\ord_Z(g)$ for $f,g\in\O_{X,Z}$.
This follows from the existence of the short exact sequence
$$ 0 \to \frac{\O_{X,Z}}{(f)} \xtos{\times g} \frac{\O_{X,Z}}{(fg)} \xtos{\proj} \frac{\O_{X,Z}}{(g)} \to 0 , $$
and the additivity of length.
This also tells us that order is a homomorphism
$$ \ord_Z\colon K(X)^\times \to \bZ . $$

Note that if $f\notin\m_{X,Z}$, meaning $f$ does not vanish along $Z$, then $f$ is invertible and so
$\O_{X,Z}/(f)=0$ and thus its order of vanishing is zero.
This is of course an equivalence.
Thus $f$ has positive order of vanishing only when it vanishes along $Z$.
And a rational function $f/g$ may have a negative order of vanishing if its denominator vanishes (``more than
its numerator'') along $Z$.

\bdefn

    Let $Y\subseteq X$ be a $(k+1)$-dimensional subvariety and let $r\in K(Y)^\times$.
    Define the $k$-cycle
    $$ \div(r) = \sum_Z\ord_Z(r)Z $$
    where the sum runs over the codimension-$1$ subvarieties of $Y$.
    Let $\rat_k(X)$ be the subgroup of $Z_k(X)$ generated by elements of the form $\sum_i\div(r_i)$
    for $r_i\in K(Y_i)^\times$ where $Y_i\subseteq X$ is a $(k+1)$-dimensional subvariety.

    Then define $X$'s {\emphcolor Chow groups} by
    $$ \Ch_k(X) = \frac{Z_k(X)}{\rat_k(X)},\qquad \Ch(X) = \bigoplus_k\Ch_k(X) . $$
    Two cycles are {\emphcolor rationally equivalent} if they have the same image in $\Ch(X)$.

\edefn

For a variety, or in general a cycle, denote its image in $\Ch(X)$ by $[\bul]$.

The Chow groups form an algebro-geometric analogue for singular homology.
Algebraic varieties over $\bC$ may be studied via the homology of their analytification, the Chow groups allow one to study varieties
over arbitrary fields.

\bexam

    If $X$ is $n$-dimensional, then its codimension-$0$ irreducible subvarieties are its irreducible components.
    As it has no $(n+1)$-dimensional subvariety, we see that $\Ch_n(X)$ is the free Abelian group generated by its irreducible components.

\eexam

Recall that rational functions are of the form $f/g$ for $f,g\in\Gamma(X,\O_X)$ and $g$ nonzero on a dense open subset of $U$.
This defines a rational function $X\to\bP^1$ on $U$ whose coordinates are $[g:f]$.
Conversely a rational function $X\to\bP^1$ is of the form $[g:f]$, and this corresponds to $f/g\in K(X)^\times$.
Thus there is a correspondence $\rat(X,\bP^1)\cong K(X)^\times$.
With this correspondence in mind we will describe an alternative characterization of rational equivalence of cycles.

Denote the distinguished points of $\bP^1$ by $0=[1:0]$ and $\infty=[0:1]$.
For a variety $V$, let $r\colon V\to\bP^1$ be a morphism, so it defines a rational function $r=f/g\in K(V)^\times$ where $r=[g:f]$.
Then the fibers of $r$ are
$$ r^{-1}0 = f^{-1}0,\qquad r^{-1}\infty = g^{-1}0 , $$
that is $r$'s zeroes are $f$'s, and its poles are $g$'s zeroes.
Let $Z_1,\dots,Z_n$ be $f^{-1}0$'s irreducible components, which are of codimension-$1$ in $X$.
Note that for $Z\subseteq X$ of codimension $1$, $\ord_Z(f)\neq0$ iff $Z=Z_i$ for some $i$.
This is because if $f\in\m_{X,Z}$ then $f$ is zero on $Z$'s generic point, and since the set of points on which $f$ is zero is closed, this means $f$
is zero on all of $Z$ (and thus zero in $\Gamma(Z,\O_X)$ as $Z$ is reduced).
And because $f=0$ on $Z$ we have that $\ord_Z(f)=\length\O_{X,Z}$.
Thus
$$ \div(r) = \sum_{Z\subseteq r^{-1}0}\length\O_{X,Z}Z - \sum_{Z\subseteq r^{-1}\infty}\length\O_{X,Z}Z , $$
where the sums are over irreducible components.

This leads us to the following definition: for a Noetherian scheme $X$ define the {\emphcolor multiplicity} of an irreducible
component $Z\subseteq X$ to be the length $m_Z=\length\O_{X,Z}$ (the ring is Noetherian of dimension zero, and thus of finite length).
Then we associate to $X$ the $\dim X$-cycle
$$ [X] = \sum_{Z\subseteq X}m_ZZ , $$
where the sum is over the irreducible components of $X$.
If $X$ is already irreducible then $\length\O_{X,X}=1$ and so $[X]$ agrees with the projection of $X\in Z_k$ into $\Ch_k$.

Then from the previous considerations,
$$ \div(r) = [r^{-1}0] - [r^{-1}\infty] . $$

\bdefn

    Let $V$ be a $(k+1)$-dimensional subvariety of $X\times\bP^1$ such that its projection onto $\bP^1$
    is dominant; denote the projection by $\pi$.
    Then for any rational point $p\in\bP^1$ the scheme-theoretic point $\pi^{-1}p$ is a $k$-dimensional
    subscheme of $X\times\set p\cong X$.
    Denote this subscheme by $V(p)\subseteq X$, so it induces a class $[V(p)]\in Z_k(X)$.

\edefn

Now it can be shown that if $V\subseteq X\times\bP^1$ is the closure of the graph of $r\colon X\to\bP^1$
then $[r^{-1}p]=[V(p)]$ and so
$$ \div(r) = [V(0)] - [V(\infty)] . $$
The converse can be shown as well (see Fulton), so we have the following.

\bprop

    A cycle $\alpha\in Z_k(X)$ is rationally equivalent to zero iff there are $(k+1)$-dimensional
    subvarieties $V_1,\dots,V_n$ of $X\times\bP^1$ such that projections $V_i\to\bP^1$ are dominant
    and in $Z_k(X)$,
    $$ \alpha = \sum_{i=1}^n[V_i(0)]-[V_i(\infty)] . $$

\eprop



Importantly the Chow group has a ring structure.
It remains graded except the grading is reversed: define $\Ch^k(X)=\Ch_{n-k}(X)$ for $n=\dim X$.

\bthrm

    For a smooth quasi-projective variety $X$, then $\Ch(X)$ (with the reverse grading) has a natural
    structure of a graded ring.
    This is called the {\emphcolor Chow ring} of $X$.

\ethrm

\bremark

    It probably pays to discuss a bit regarding how the ring structure is defined.
    The naive structure
    $$ [A][B] = [A\cap B] $$
    for subvarieties $A,B$ is not well-defined.
    It is well-defined when we place some conditions on $A,B$, specifically that they are
    {\emphcolor generically transverse}.
    Recall that $A$ and $B$ intersect transversally at $p\in A\cap B$ if $A,B,X$ are smooth at $p$ and
    $T_pA+T_pB=T_pX$.
    Then $A$ and $B$ are generically transverse if there is a dense subset of their intersection over which
    they intersect transversally.
    Two cycles $A=\sum_in_iA_i$ and $B=\sum_jm_jB_j$ are generically transverse if all $A_i$ and $B_j$ are.

    Then it can be proven with much effort that for every $\alpha,\beta\in\Ch(X)$ there are generically
    transverse representatives $A,B$ of $\alpha,\beta$.
    Furthermore $[A\cap B]$ is independent on choice of generically transverse cycles.
    This gives the ring structure.

\eremark

